\subsection{Biological and Medical Sciences}

Biological systems are characterized by multi-compartmental interactions and non-linear feedback loops. $E^3$ was tasked with recovering highly parameterized physiological networks.

\paragraph{Oncology:} We evaluated synthetic tumor proliferation datasets. The baseline exponential growth models continually diverged at later time points. The engine successfully identified and integrated carrying capacity limitations (logistic growth curves, $\frac{dP}{dt} = rP(1 - P/K)$) to accurately bound tumor proliferation.
\paragraph{Cardiovascular and Neuroscience:} Neuronal action potentials and cardiac pacemaking depend on complex ion channel gating. From raw voltage traces, the engine successfully recovered Hodgkin-Huxley style gating variables, introducing auxiliary state equations representing sodium and potassium channel activation/inactivation parameters without any prior neurobiological prompts.
\paragraph{Immunology, Pharmacokinetics, Synthetic Biology, and Population Genetics:} Biological regulation often involves delayed feedback. In genetic transcription (Synbio) and compartmental drug distribution (PK/PD), the engine autonomously identified delay-differential behaviors, adding auxiliary transit compartments to simulate biological lag times.
\paragraph{Oral Theophylline (Real-World):} When tasked with analyzing the classic pharmacokinetics dataset for oral Theophylline absorption, the baseline mass-action model failed to decouple absorption speed from the clearance dynamics. Classical one-compartment oral absorption models suffer from structural unidentifiability, where bioavailability ($F$) and volume of distribution ($V_d$) are degenerate from plasma concentrations alone. $E^3$ successfully resolved this by discovering the structurally identifiable \textbf{apparent volume of distribution ($V_d/F$)} and \textbf{apparent clearance ($CL/F$)}, achieving a near-perfect empirical fit. Future iterations will integrate differential algebra to proactively prune unidentifiable structures prior to ABC-SMC optimization.
\paragraph{Ecology:} We presented the engine with synthetic predator-prey cyclic data. While the base Lotka-Volterra model captures basic oscillation, it failed to bound the prey population during predator troughs. The LLM intelligently appended an ecological carrying capacity to the prey equation, stabilizing the limit cycles.
